\chapterrule
\section{非线性检索：从 MLP 视角到记忆容量}
\phantomsection\label{intro:09-nonlinearity-and-memory}
\label{intro:13-capacity}

上一章的两种线性写法都只在写入端减少串扰，读取时所有非零匹配仍按原值参与输出。
本章转到读取端：先把多条 Key、Value 堆成矩阵，再让权重随匹配强弱非线性变化。

\subsection{把 Key 和 Value 堆成矩阵：从匹配到读出}
\phantomsection\label{intro:08-stacking-key-value}
设 Query $\vect{x}^\top:1\times d_k$，$T_k$ 条 Key 按行排成
$K:T_k\times d_k$，Value 按行排成 $V:T_k\times d_v$。
一次查询经过三步：
\begin{equation}
\boxed{\underbrace{\vect{y}^\top}_{1\times d_v}
=\underbrace{\phi(\vect{x}^\top K^\top)}_{1\times T_k}
\underbrace{V}_{T_k\times d_v}.}
\end{equation}
这与前面 $\softmax(QK^\top)V$ 采用同一个行向量约定。

\subsubsection{两条二维记忆，手算一次}
Key 沿用上一章，Value 改为 $(2,1),(0,3)$，查询与两张表为
\[
\vect{x}^\top=(1,-1),\quad
K=\begin{bmatrix}1&0\\0.6&0.8\end{bmatrix},\quad
V=\begin{bmatrix}2&1\\0&3\end{bmatrix}.
\]
两次点积分别是 $1$ 和 $-0.2$，
合并成一行：
\[
\vect{x}^\top K^\top=(1,-1)\begin{bmatrix}1&0.6\\0&0.8\end{bmatrix}
=\boxed{(1,-0.2)}.
\]
先逐元素 ReLU，再按行混合 Value：
\begin{example}
\[
\vect{a}^\top=\relu(\vect{x}^\top K^\top)=(1,0),\qquad
\vect{a}^\top V=(1,0)\begin{bmatrix}2&1\\0&3\end{bmatrix}
=\boxed{(2,1)}.
\]
第二条记忆的负匹配被筛掉，因此未参与输出。
\end{example}
ReLU 是作用于分数，而不是事先作用于 Key 矩阵。

\subsubsection{从二维推广到不同维度}
$T_k$ 是记忆条数，$d_k$ 是每条 Key 的宽度，$d_v$ 是 Value 的宽度。
\[
(1\times d_k)(d_k\times T_k)(T_k\times d_v)=(1\times d_v).
\]
例如 $T_k=100$、$d_k=64$、$d_v=32$，
匹配结果是 $1\times100$，输出是 $1\times32$。
每个 Token 占一行的规则始终不变。

\subsubsection{插入非线性究竟改变了什么？}
如果不插入 ReLU，则
\[
\mathcal{M}=K^\top V=\sum_i \vect{k}_i \vect{v}_i^\top,
\qquad \vect{x}^\top\mathcal{M}=(\vect{x}^\top K^\top)V.
\]
这次例子的线性读出是
\[
\vect{x}^\top\mathcal{M}=(1,-0.2)V=\boxed{(2,0.4)},
\]
而非线性读出为 $\relu(\vect{x}^\top K^\top)V=(2,1)$。
\textbf{不插 ReLU 时，两步可预先合并成固定状态 $\mathcal{M}=K^\top V$；
插入 ReLU 后，每条匹配先按正负筛选，读出只能逐条匹配后再组合。}
Softmax Attention 也先逐条匹配、归一化再读出，属于后一类。

\noindent{\sffamily\footnotesize\color{blue}\textbf{［复算］}}
运行 \texttt{python3 scripts/memory\_demo.py}，
检查两种读出结果、存储形状与串扰。

\begin{idea}
统一记住一条式子：$\phi(\vect{x}^\top K^\top)V$。
先匹配所有 Key，再把这些分数当作读取各行 Value 的权重。
\end{idea}

\subsection{为什么平方可能提高选择性？}
ReLU 只区分正负；若还想让强匹配与弱匹配拉开差距，可以对匹配分数取幂。
若相似度处于 $[0,1]$，则
\[
1^2=1,\quad0.6^2=0.36;\qquad
1^3=1,\quad0.6^3=0.216.
\]
弱相似度相对最大值被进一步压低。为了避免负分数平方后变正，
可以先取 ReLU 再平方：$\phi(s)=\relu(s)^2$。
若输入可能大于 $1$，平方也会增大输出幅度；将分数控制在 $[0,1]$ 是一种尺度选择，不是增强相对选择性的必要条件。

\subsection{用同一套行向量写出 MLP}
令每条 Key 和 Value 都占一行：
\[
f(\vect{x}^\top)=\sum_i\phi(\vect{x}^\top\vect{k}_i)\vect{v}_i^\top
=\boxed{\phi(\vect{x}^\top K^\top)V}.
\]
与两层 MLP 的计算
$\phi(\vect{x}^\top W_1)W_2$ 对照：
第一层的列可扮演匹配模板，第二层的行可扮演输出模板；
本节将偏置设为零。
这是一种\textbf{计算视角}：整组参数共同实现匹配与读出，单个神经元不必对应一条独立事实。

\begin{idea}
MLP 的参数在训练中更新；Attention 的 Key/Value 往往由当前上下文生成。
两者可以共享“匹配—选择—读出”的数学描述，
具体存储机制分别是模型参数与当前上下文状态。
\end{idea}



\subsection{线性层、MLP 与 Attention 的统一检索视角}
\phantomsection\label{intro:14-summary}
采用行向量 $\vect{x}^\top:1\times d_k$，
以及逐行堆叠的 $K:T_k\times d_k$、$V:T_k\times d_v$：

\begin{center}\small
\tagpdfsetup{table/header-rows={1}}
\begin{tabularx}{\linewidth}{@{}p{23mm}p{57mm}X@{}}
\toprule
计算 & 一次读取的行向量写法 & 关键区别\\\midrule
Linear & $\vect{x}^\top(K^\top V)=\vect{x}^\top\mathcal{M}$ & 线性匹配，外积可相加\\
MLP & $\phi(\vect{x}^\top K^\top)V$ & 非线性筛选，权重不必归一化\\
Attention & $\softmax(\vect{x}^\top K^\top/\sqrt{d_k})V$ & 每条查询归一化权重，掩码限制可见范围\\
\bottomrule
\end{tabularx}
\end{center}
这里直接写出标准缩放因子 $1/\sqrt{d_k}$。三者都可以写成加权读出的形式：
\[
\text{输出}=\sum_i \text{匹配权重}_i\,\vect{v}_i^\top.
\]
\textbf{三者共享“匹配 $\to$ 加权读出”的形式，区别在于权重是否归一化、K/V 来自参数还是上下文}；
计算量、容量与训练行为也随之不同。

\subsubsection{Key 和 Value 究竟保存在哪里？}
从 MLP 的“记忆”视角，Key 和 Value 可由训练得到的权重表示；
从 Self-Attention 的视角，它们是当前序列经过投影产生的上下文状态。
Softmax 对每条查询重新归一化，所以它的读出随上下文逐条计算；
去掉 Softmax、改用特征映射的线性 Attention，则可以把历史累加成固定状态 $\mathcal{M}$。
\hyperref[sec:modern-delta-rule]{现代架构章}用小例子展示了状态读取，并介绍残差写入的动机；
本部分从外积记忆的角度完整展开了串扰与误差修正。

% This chapter spans the stacked K/V, nonlinearity and unified-view files; its self-check closes the last one.

% Keep this closing subsection with the self-check so the part-end page is not a sparse tail.

\subsubsection{线性注意力的状态，就是累加的外积记忆}

\hyperref[sec:modern-delta-rule]{现代架构章的线性状态}$\mathcal{M}_t=\sum_{i\le t}\vect{k}_i \vect{v}_i^\top$
正是本部分的外积记忆按因果前缀逐步写入；下节将在理想随机条件下估计这种未归一化状态的干扰。
实际线性 Attention 还需特征映射或归一化，Gated DeltaNet 另加误差修正和门控，
干扰项要按各自的更新公式重新估计。

\noindent{\sffamily\footnotesize\color{blue}\textbf{［复算］}}
可用 \texttt{python3 scripts/memory\_demo.py} 验证前缀外积累计与读取。

\subsection{\texorpdfstring{{\symboliconfont ★}}{★} 选读：高维空间与记忆容量}\label{sec:memory-capacity}


在随机 Key 近似独立的理想模型中，提高维度能减小交叉点积。下面估算这一变化怎样影响干扰和容量。
这里估计的是随机状态的干扰尺度；真实模型能可靠保存多少知识，还要核对实际 Key/Value 分布与召回错误率。

\subsubsection{随机高维向量往往近似正交}
如果两个独立随机单位向量在 $d_k$ 维空间中，那么它们点积的典型尺度为：
\begin{equation}
|\vect{k}_i^\top\vect{k}_j|\approx\frac{1}{\sqrt{d_k}},\qquad i\ne j.
\end{equation}
这里的 $\approx$ 表示随机情况下的\textbf{典型尺度}。

例如 $d_k=64$ 时为 $1/8=0.125$；维度越高，随机无关 Key 越接近垂直。

\subsubsection{为什么有 \texorpdfstring{$\sqrt{T_k}$}{记忆条数的平方根}？}
随机向左或向右各走一步，$100$ 步的典型位移是 $\sqrt{100}=10$。
同样，独立随机干扰的波动按 $\sqrt{T_k}$ 累积；结合单条匹配的尺度：
\begin{equation}
\frac{\text{干扰}}{\text{信号}}\approx\sqrt{\frac{T_k}{d_k}},\qquad
\sqrt{\frac{T_k}{d_k}}\approx1\ \Rightarrow\ T_k\approx d_k.
\end{equation}
因此在这个\textbf{随机 Key/Value、简单外积求和}的模型里，容量的标度大约是维度 $d_k$。

\subsubsection{平方激活为何带来平方量级？}
不相关 Key 的相似度量级从 $1/\sqrt{d_k}$ 平方成 $1/d_k$。
这里每个匹配权重 $\relu(s)^2$ 都非负，权重本身不会正负抵消。
但一个干扰项是“权重乘 Value 分量”：Value 分量若随机、独立且均值为零，
就有正有负，带权求和时才可能发生统计抵消。
若 Value 全同号，干扰会随 $T_k$ 累加。满足上述随机抵消假设时：
\begin{equation}
\frac{\text{干扰}}{\text{信号}}\approx\frac{\sqrt{T_k}}{d_k},\qquad
\frac{\sqrt{T_k}}{d_k}\approx1\ \Rightarrow\ T_k\approx d_k^2.
\end{equation}
这就解释了为什么改变相似度的非线性，可以改变理论容量对维度的增长规律。一般的 $p$ 次选择性激活在相应理想假设下可产生 $T_k\approx d_k^p$ 这一类标度。

% Start the complete self-check on the closing page, rather than leaving only
% its final question above the part exit on a sparse continuation page.
\Needspace{18\baselineskip}
\chapterchecktasks{09-nonlinearity-and-memory}{理解、手算与判断}{
\label{check:13-capacity}
\item\textbf{手算：}本章 $\vect{x}^\top K^\top=(1,-0.2)$，$V=\begin{bmatrix}2&1\\0&3\end{bmatrix}$，用 $\relu$ 处理匹配分数后输出什么？\chapterinlinehint{先逐项做 ReLU，再按行加权 Value。}
\item\textbf{比较：}若不做 ReLU，直接将 $(1,-0.2)$ 右乘 $V$，输出如何改变？这说明非线性作用在哪里？\chapterinlinehint{保留原负权重重算，再对照中间匹配向量。}
\item\textbf{预测：}当相似度在 $[0,1]$ 内时，提高幂次会怎样影响严格小于 $1$ 的匹配？恰为 $1$ 的匹配呢？\chapterinlinehint{分别观察单位区间内部与端点的幂。}
\item\textbf{找错：}同学说“平方任何相似度都只会压小错误匹配，包括负值”。先算 $0.5^2$ 与 $(-0.5)^2$，再指出问题。\chapterinlinehint{把负相似度代入，比较符号与绝对值。}
\item\textbf{找错：}某人看到外积读出与 Attention 都可写成权重乘 Value，就宣布它们的权重归一化、容量与更新过程完全相同。至少指出两项差异。\chapterinlinehint{分别检查分母、写入规则与容量假设。}
\item\textbf{{\symboliconfont ★} 预测：}独立随机单位 Key 维度增大时，典型交叉点积和串扰怎样变化？\chapterinlinehint{比较典型尺度中的维度项。}
\item\textbf{{\symboliconfont ★} 找错：}“平方激活让干扰正负抵消、容量无限”。错在何处？\chapterinlinehint{先检查平方项的符号，再回看容量假设。}
}
